\documentclass[12pt]{article}
\usepackage[T1]{fontenc}
\usepackage[utf8]{inputenc}
\usepackage{lmodern}
\usepackage{textcomp}
\usepackage{enumitem}
\usepackage{geometry}
\geometry{margin=1in}
\begin{document}
\begin{center}
Answer Key - Chemistry - Undergraduate Level Assessment 3 \
\small (Answers are ordered exactly as in the question paper.)
\end{center}

\section*{Multiple Choice}
\begin{enumerate}[label=\arabic*.]
\item \textbf{Answer:} A
\\ \textit{Options:} A) PbH 4 \textless{} SnH 4 \textless{} GeH 4 \textless{} SiH 4 \textless{} CH 4; B) PbH 4 \textless{} SnH 4 \textless{} CH 4 \textless{} GeH 4 \textless{} SiH 4; C) CH 4 \textless{} SiH 4 \textless{} GeH 4 \textless{} SnH 4 \textless{} PbH 4; D) CH 4 \textless{} PbH 4 \textless{} GeH 4 \textless{} SnH 4 \textless{} SiH 4
\\ \textit{Note:} The thermal stability of hydrides of group-14 elements increases with decreasing atomic size and increasing bond strength, making CH 4 the most stable and PbH 4 the least stable due to the weakening of the M-H bond as the size of the central atom increases.
\item \textbf{Answer:} D
\\ \textit{Options:} A) 1.5 R; B) 2 R; C) 2.5 R; D) 3.5 R
\\ \textit{Note:} The limiting high-temperature molar heat capacity at constant volume (C\_V) of a gas-phase diatomic molecule is 3.5 R because, at high temperatures, the molecule can access three translational and two rotational degrees of freedom, contributing 5/2 R from translational motion and 2/2 R from rotational motion, resulting in a total of 3.5 R.
\item \textbf{Answer:} B
\\ \textit{Options:} A) 0.0471 Hz; B) 21.2 Hz; C) 42.4 Hz; D) 66.7 Hz
\\ \textit{Note:} The linewidth can be calculated using the formula Δν = 1/(2πT 2), and substituting T$_{2}$ = 15 ms yields a linewidth of approximately 21.2 Hz.
\item \textbf{Answer:} D
\\ \textit{Options:} A) I only; B) II only; C) III only; D) I and III only
\\ \textit{Note:} The normal modes of a carbon dioxide molecule that are infrared-active are those that result in a change in the dipole moment, which occurs in the bending mode and the asymmetric stretching mode, but not in the symmetric stretching mode, as it does not produce a net dipole change.
\item \textbf{Answer:} D
\\ \textit{Options:} A) Na$_{2}$CO$_{3}$; B) Na$_{3}$PO$_{4}$; C) Na$_{2}$S; D) NaCl
\\ \textit{Note:} The correct answer is NaCl because it is a neutral salt formed from a strong acid (HCl) and a strong base (NaOH), resulting in a solution that does not significantly affect the pH, while other options may have acidic or basic properties that could lower or raise the pH.
\end{enumerate}

\section*{True / False}
\begin{enumerate}[label=\arabic*.]
\setcounter{enumi}{5}
\item \textbf{Answer:} True
\\ \textit{Options:} A) True; B) False
\\ \textit{Note:} Protons in a molecule with a rotational correlation time of 1 ns experience optimal conditions for spin-lattice relaxation at a magnetic field of 3.74 T due to the balance between the molecular motion and the magnetic environment, which enhances energy transfer processes.
\item \textbf{Answer:} False
\\ \textit{Options:} A) True; B) False
\\ \textit{Note:} The statement is false because the NMR frequency of 19 F in a 20.0 T magnetic field is approximately 79.1 MHz, not 115.5 MHz.
\item \textbf{Answer:} False
\\ \textit{Options:} A) True; B) False
\\ \textit{Note:} The statement is false because AgCl, being an ionic compound, actually has a higher melting point than HCl, which is a molecular compound with weaker van der Waals forces.
\item \textbf{Answer:} True
\\ \textit{Options:} A) True; B) False
\\ \textit{Note:} Most organic compounds are primarily composed of carbon, hydrogen, oxygen, and nitrogen, which means that the presence of silicon is rare, making tetramethylsilane, which contains no hydrogen atoms on adjacent carbons and serves as a non-polar solvent, an ideal reference for 1 H NMR spectroscopy.
\item \textbf{Answer:} False
\\ \textit{Options:} A) True; B) False
\\ \textit{Note:} The statement is false because thionyl chloride, SOCl 2, has three bonding pairs of electrons and one lone pair, resulting in a trigonal pyramidal molecular geometry rather than a tetrahedral shape.
\end{enumerate}

\section*{Long Form Response}
\begin{enumerate}[label=\arabic*.]
\setcounter{enumi}{10}
\item \textbf{Answer:} The strength of a 13 C$_{90}$^\ensuremath{\circ} pulse with a duration of 1 μs is influenced by several factors, including the gyromagnetic ratio of 13 C, the pulse duration, and the desired flip angle. For 13 C, the gyromagnetic ratio is approximately 67.28 MHz/T. To achieve a 90^\ensuremath{\circ} flip angle, the pulse strength (B$_{1}$) can be calculated using the formula B$_{1}$ = (π/2) / (γ \ensuremath{\times} τ), where τ is the duration of the pulse. Substituting the values, we find B$_{1}$ \ensuremath{\approx} 5.18 mT. This value ensures that the magnetic field strength is sufficient to induce the desired nuclear spin transition in 13 C nuclei during the specified pulse duration.
\\ \textit{Note:} The answer correctly identifies that the strength of a 13 C$_{90}$° pulse is determined by the gyromagnetic ratio and pulse duration, using the appropriate formula to calculate the required magnetic field strength to achieve the desired flip angle.
\item \textbf{Answer:} Dissociation energy is the energy required to break a bond in a molecule, resulting in the formation of individual atoms. In the case of the reaction HBr(g) → H(g) + Br(g), the dissociation energy corresponds to the enthalpy change associated with breaking the H-Br bond in hydrogen bromide. This process is endothermic, meaning it requires an input of energy to overcome the attractive forces between the hydrogen and bromine atoms. The dissociation energy for the H-Br bond is approximately 366 kJ/mol, which reflects the stability of the HBr molecule and the energy needed to separate it into its constituent atoms. Understanding this concept is crucial for evaluating reaction mechanisms and the energy profiles of chemical processes.
\\ \textit{Note:} The answer correctly identifies dissociation energy as the energy needed to break the H-Br bond in hydrogen bromide, characterizing the process as endothermic and providing the specific value of approximately 366 kJ for the dissociation energy, which exemplifies the concept in the context of the reaction.
\end{enumerate}

\vfill
\noindent\textit{End of Answer Key}
\end{document}